% =============================================================
%  ONLINE RESOURCE 4 — Architecture Ablation
%  Supporting Information for Hackman & Zhu (2026)
%  "When Inverse Physics-Informed Neural Networks Become
%   Practically Identifiable: A Case Study in Synaptic Dopamine
%   Transport" — Medical & Biological Engineering & Computing
%   submission
%
%  Note: sections A-D of this resource (forward-problem
%  space-time error field, noise sensitivity, parameter grid,
%  observation density scaling) were restored to the main
%  manuscript as Figs. 3, 5, 7, 8 to match the figure set used
%  in the PLOS Computational Biology submission. Only the
%  architecture ablation (never a PLOS main figure) remains here.
% =============================================================
\documentclass[12pt]{article}

\usepackage[margin=1in]{geometry}
\usepackage[utf8]{inputenc}
\usepackage[T1]{fontenc}
\usepackage{lmodern}
\usepackage{amsmath, amssymb, mathtools}
\usepackage{graphicx}
\usepackage{booktabs}
\usepackage{hyperref}
\hypersetup{colorlinks=true, linkcolor=black, citecolor=black,
            urlcolor=blue, pdfborder={0 0 0}}
\usepackage{setspace}
\setlength{\parindent}{0pt}
\setlength{\parskip}{0.7em}

\title{Online Resource 4: Architecture Ablation \\
\large Supporting Information for ``When Inverse
Physics-Informed Neural Networks Become Practically
Identifiable: A Case Study in Synaptic Dopamine Transport''}

\author{Emmanuel Hackman and Huiqing Zhu \\
School of Mathematics and Natural Sciences, \\
University of Southern Mississippi, Hattiesburg, MS 39406, USA}

\date{}

\begin{document}
\maketitle

This resource reports the architecture ablation study
referenced in the Results section of the main text.

\section*{Architecture Ablation}

To verify that the canonical four-hidden-layer, 64-neuron
architecture is not badly over- or under-parameterized, we
re-trained the forward problem with three MLP configurations on
identical data and hyperparameters.

\begin{table}[h!]
\centering
\caption{\textbf{Architecture ablation.}
Forward-problem $L^2$ error against the finite-difference reference and
total training time (Adam $+$ L-BFGS, single seed $1234$) for three
MLP configurations.}
\label{tab:arch_ablation}
\begin{tabular}{lccc}
\toprule
\textbf{Configuration}                 & \textbf{Params} & \textbf{$L^2$ vs FD (\%)} & \textbf{Train time (s)} \\
\midrule
$[3, 32, 32, 32, 1]$ (compact)         & $2{,}273$  & $1.38$ & $241.1$ \\
$[3, 64, 64, 64, 64, 1]$ (canonical)   & $12{,}801$ & $1.52$ & $329.5$ \\
$[3, 128, 128, 128, 1]$ (wide)         & $33{,}665$ & $2.04$ & $308.6$ \\
\bottomrule
\end{tabular}
\end{table}

All three configurations achieve sub-$2.1\%$ forward $L^2$ error;
the wide configuration is slightly less accurate ($2.04\%$),
likely reflecting mild overparameterization for this smooth
2D problem. The canonical architecture is competitive with both
alternatives and robust to moderate changes in capacity.

\end{document}
